% !PRINTED 508
\weberpage{508}
%<<CONT
wenn $a\partial = n$ und $c$ durch 8 teilbar ist, eine ganze algebraische
Funktion des Körpers $\mathfrak{R}(v,u)$, worin
\weq{2}{v = j\left(\frac{c+\partial\omega}{2}\right), \qquad u = j(\omega).}
% [sic] printed denominator $2$ conflicts with denominator $a$ in (5) and with the ensuing system (7)
Es werde nun darin für $\omega$ eine Wurzel der primitiven
quadratischen Gleichung
\weq{3}{A\omega^2 + B\omega + C = 0,}
mit negativer Diskriminante
\weq{4}{D = B^2 - 4AC.}
% [sic] printed full stop interrupts the sentence, whose governing verb ``gesetzt'' follows the display
gesetzt, und wir nehmen ein für allemal an, daß $A$ positiv und
relativ prim zu $2Dn$ sei, was keine Beschränkung ist. Soll nun
$v = u$ werden, so ist dafür die notwendige und hinreichende
Bedingung:
\weq{5}{\frac{c+\partial\omega}{a} = \frac{\gamma+\delta\omega}{\alpha+\beta\omega},}
\weq{6}{\alpha\delta - \beta\gamma = 1.}
Vergleicht man (5) mit (3), so folgt (§ 114), daß es zwei
ganze positive Zahlen $x, y$ geben muß, deren erste positiv angenommen werden kann und die den Bedingungen genügen:
\weq{7}{\begin{aligned}
\partial\beta &= Ax, & c\beta+\partial\alpha-a\delta &= Bx,\\
c\alpha-a\gamma &= Cx, & -c\beta+\partial\alpha+a\delta &= y,\\
2\partial\alpha &= Bx+y, & & \\
2(c\beta-a\delta) &= Bx-y, & &
\end{aligned}}
und daraus wegen (6):
\weq{8}{4n = y^2 - Dx^2.}
Ein Primteiler von $\partial$ müßte, da $A$ relativ prim zu $n$ angenommen ist, in $x$ und in $y$ aufgehen. Dann aber auch in
\[
c\alpha - a\gamma \quad\text{und in}\quad c\beta - a\delta,
\]
folglich in $a, \partial, c$. Da diese drei Zahlen aber ohne gemeinschaftlichen Teiler sind, so muß $\partial = 1$ sein, und aus (7) ergibt sich:
\weq{9}{\begin{aligned}
\alpha &= \frac{Bx+y}{2}, & \beta &= Ax,\\
c\alpha - n\gamma &= Cx, & c\beta - n\delta &= \frac{Bx-y}{2};
\end{aligned}}
daraus folgt, daß $x$ und $y$ keinen ungeraden gemeinschaftlichen Teiler haben, und daß, wenn $x$ und $y$ gerade sind, $\frac{1}{2}(Bx+y)$ ungerade sein muß. Da wir überdies $x$ positiv annehmen können,
%CONT>>
